\documentclass{article}

\usepackage{xcolor}
\usepackage{tikz}
\usetikzlibrary{trees}
\usepackage[utf8]{inputenc}
\usepackage{amsmath}
\usepackage{enumerate}
\usepackage{microtype}
\usepackage[top=1.5cm]{geometry} % controls page margins

\title{Examen Parcial 2}
\author{Programación III \\ Facultad de Ingenierías - Universidad Tecnológica de Pereira \\ Ingeniería de Sistemas y Computación}
\date{Abril 21 2026}
\begin{document}
\maketitle
\thispagestyle{empty}

Un estudiante cuenta con un ingreso semanal fijo de \$400.000, a partir del cual debe planear su rutina y gasto total para un espacio de 7 días.

Considere las siguientes reglas que se deben satsifacer en una semana de estudio:

\begin{itemize}
  \item Antes del tercer día de la semana debe pagar una parte del arriendo de su apartamento, con un costo de \$200.000. Una vez ha pagado su arriendo, no debe volver a preocuparse por ello.
  \item Cada día, el estudiante debe elegir una de las siguientes opciones de comida:
  \begin{itemize}
    \item Opción económica (\$5.000): Bajo costo, insatisfactoria.
    \item Opción intermedia(\$10.000): Costo medio, un poco satisfactoria.
    \item Opción costosa(\$25.000): Costo alto, muy satisfactoria.
  \end{itemize}
  Sin embargo, debido a que la opción económica no es para nada saludable, el estudiante no puede elegirla más de 2 veces a la semana, y no puede ser en días contiguos.
  
  (\textbf{Opcional(+0.7)} El estudiante es una persona que no se siente satisfecho con la monotonía, por lo tanto no está dispuesto a elegir el mismo tipo de comida tres días seguidos)
  \item Durante la semana, el estudiante debe lavar su ropa en la lavandería; cada visita tiene un costo fijo de \$ 20.000:
  \begin{itemize}
    \item En una semana, debe visitar la lavandería por lo menos dos veces a la semana
    \item Estas visitas no pueden ocurrir en dos días consecutivos, así como tampoco el mismo día en el que se pagó el arriendo.
  \end{itemize}
  \item Para mantenerse al día con sus lecciones, el estudiante debe revisar el material de estudio de sus clases por lo menos una vez cada dos días.
  \item Para mantenerse en contacto con sus amigos, el estudiante debe dedicar por lo menos uno de los días de su semana a una actividad de esparcimiento; para lo cual tiene dos opciones: 
  \begin{itemize}
    \item Una actividad de esparcimiento sencilla (\$30.000) Es algo satisfactoria.
    \item Una actividad de esparcimiento elaborada (\$60.000) Es muy satisfactoria.
  \end{itemize}
\end{itemize}

Debido a que cuenta con tiempo limitado, el estudiante sólo puede realizar una acción al día, además de comer, la cual puede ser:
\begin{itemize}
  \item Pagar el arriendo (1 vez).
  \item Estudiar
  \item Ir a la lavandería.
  \item Esparcimiento.
\end{itemize}

De manera adicional, se espera que el estudiante logre terminar su semana sintiéndose tan satisfecho como sea posible y habiendo ahorrado la mayor cantidad de dinero posible. Considere que el estudiante no puede gastar en la semana una suma de dinero mayor a su ingreso semanal.
\end{document}